\documentclass[11pt,a4paper]{article}
\usepackage[margin=25mm]{geometry}
\usepackage{amsmath,amssymb,amsthm,hyperref}
\setlength{\emergencystretch}{2em}
\newtheorem{lemma}{Lemma}\newtheorem{theorem}[lemma]{Theorem}
\newtheorem{corollary}[lemma]{Corollary}
\newcommand{\gen}[1]{\langle#1\rangle}
\newcommand{\PSL}{\operatorname{PSL}}
\newcommand{\Ad}{\operatorname{Ad}}
\title{Completing the odd projective-linear triality geometries}
\author{Human Compression project: L023 uniform completion proof}\date{}
\begin{document}\maketitle
The unrestricted Construction 6 in Leemans--Stokes,
arXiv:2208.08215v2, need not produce a geometry, even from a faithful
Class III map. This paper proves the missing uniform completion for
the odd-characteristic projective-linear family. It therefore supplies
actual infinitely many geometries, rather than extrapolating a finite
test. The source's advertised family belongs to its authors; this
completion proof repairs the missing all-flags step. All programs here
are optional checks of the written algebra, not finite premises.

\section{A signed-vector intersection bound}
Let $K$ be a finite field of odd characteristic. Distinct commuting
involutions $a,c\in\PSL_2(K)$ have lifts $i,j\in\operatorname{SL}_2(K)$
with $i^2=j^2=-I$ and $ij=-ji$. Indeed a noncentral determinant-one
lift of a projective involution is trace zero. If its lift commuted
with another such lift, the two trace-zero matrices would be
proportional in its two-dimensional matrix centralizer, giving the
same projective involution. Thus distinct projectively commuting
involutions have anticommuting lifts. Put $k=ij$.

The matrices $i,j,k$ form an orthonormal basis for the trace-zero
space $W$, with bilinear form
\[
 (X,Y)=-\tfrac12\operatorname{tr}(XY),\qquad
 N(X)=(X,X),\qquad X^2=-N(X)I.
\]
Their multiplication is the quaternion multiplication
$ij=k$, $jk=i$, $ki=j$, with reversed products negative.
Conjugation gives a faithful action of $\PSL_2(K)$ on $W$,
preserving this form and with determinant one. The kernel consists
of scalar matrices. For the determinant statement, conjugation on
$M_2(K)$ has determinant $(\det A)^2(\det A^{-1})^2=1$;
it fixes the scalar summand and preserves the trace-zero summand.

Take a third involution $b$ with lift
\[
                 B=xi+yj+zk,\qquad x^2+y^2+z^2=1.
\]
Suppose $b$ commutes with none of $a,c,ac$. Then $x,y,z$ are
all nonzero. For example $x=0$ would make $B$ anticommute with $i$,
and hence make $b$ commute projectively with $a$. Define
\[
 v_i=(0,-z,y),\qquad v_j=(z,0,-x),\qquad v_k=(-y,x,0)
\]
in the basis $i,j,k$. These are the commutators with $B$, divided
by two. Each belongs to $B^\perp$ and is nonzero. Any two are
independent: for example $v_i\times v_k=-yB\ne0$.
They consequently span $B^\perp$, which is nondegenerate because
$N(B)=1$. Their pairwise inner products are
\[
 (v_i,v_j)=-xy,\quad(v_j,v_k)=-yz,\quad(v_k,v_i)=-zx,
\]
all nonzero. Individual $v_i$ may be isotropic; we never divide
by their norms.

For a norm-one vector $B$, direct multiplication gives
\[
                  \Ad_b(X)=2(B,X)B-X.                 \tag{1}
\]
It fixes $B$ and negates $B^\perp$. The involutions $a$ and $b$
therefore both negate $v_i$; similarly $c,b$ negate $v_j$, and
$ac,b$ negate $v_k$. Thus
\[
 H_i=\gen{a,b},\quad H_j=\gen{c,b},\quad H_k=\gen{ac,b}
\]
respectively stabilize the signed vectors $\{v_i,-v_i\}$,
$\{v_j,-v_j\}$, $\{v_k,-v_k\}$.

\begin{lemma}\label{stabilizer}
In the adjoint group, the simultaneous stabilizer of any two of
these signed vectors is exactly $\{1,b\}$. Consequently
$H_i\cap H_j=H_i\cap H_k=H_j\cap H_k=\gen b$.
\end{lemma}
\begin{proof}
An isometry sending two of the vectors to signed copies must use
the same sign on both, since their inner product is nonzero. It is
therefore $+1$ or $-1$ on all of $B^\perp$. It preserves the
orthogonal line $KB$; its determinant one forces its action on that
line to be $+1$. Thus it is the identity or (1). Faithfulness of the
adjoint action identifies these transformations with $1,b$.
Both lie in each displayed parabolic, proving the intersections.
\end{proof}

\begin{lemma}\label{three}
There are at most three signed vectors $\{v,-v\}$ satisfying
\[
 N(v)=N(v_j),\quad (v,v_i)=\pm(v_j,v_i),\quad
 (v,v_k)=\pm(v_j,v_k).
\]
\end{lemma}
\begin{proof}
Choose the unique sign of $v$ making $(v,v_i)=(v_j,v_i)$;
the latter is nonzero. If also $(v,v_k)=(v_j,v_k)$, the two
independent linear equations give $v=v_j+tB$. Because
$(v_j,B)=0$ and $N(B)=1$, equality of norms forces $t^2=0$,
and hence $v=v_j$. For the other sign of the second inner product,
the solutions of the two linear equations form an affine line
parallel to $B$. Restricting the norm equation to this line gives
a polynomial in $t$ of degree two with leading coefficient one.
There are at most two further solutions. This argument also applies
when $N(v_j)=0$.
\end{proof}

\section{All flags complete, with the exact residue}
Put $G=\gen{a,b,c}$, $V=\gen{a,c}$, and use the four subgroups
$H_i,V,H_j,H_k$ for faces, edges, vertices and Petrie objects.
The intersection of $H_i$ with $V$ is $\gen a$, and the other
two intersections are $\gen c$ and $\gen{ac}$. For example
$\Ad_c(v_i)=(0,-z,-y)$ is neither $v_i$ nor $-v_i$, because
$y,z\ne0$; the same applies to $ac$, and the other cases follow
cyclically. In particular the three nonedge parabolics are proper.

The rank-three truncation is a faithful thin surface map. At each
edge its two faces and two vertices give four flags, a free orbit
of $V$. Every incident face--vertex pair extends using an actual
element of the coset intersection; Lemma~\ref{stabilizer} shows it
has exactly two edges. Thus the flags are a single regular $G$-orbit,
and all panels are thin. The flag-adjacency graph is connected by
generation. Each vertex or face residue is the cycle of its marked
dihedral parabolic. Gluing the corresponding flag triangles therefore
gives a closed surface with a two-cell embedded graph. Conversely,
unique flag adjacencies imply that a map automorphism fixing a flag
fixes all flags. Its full type-preserving incidence group is exactly $G$.

\begin{theorem}\label{completion}
For any such odd-characteristic marking, the four-type coset
pregeometry is an incidence geometry. The residue of an incident
face--Petrie pair is precisely a path with three vertices and two
edges. There are two free $G$-orbits of chambers, and the geometry
is not firm.
\end{theorem}
\begin{proof}
It suffices to consider the standard pair $H_i,H_k$, since an
incident pair of cosets has a common element and can be translated
to this pair. A vertex $gH_j$ incident with both has a well-defined
signed vector $\{\Ad_g(v_j),-\Ad_g(v_j)\}$. Choose a representative
of this coset in $H_i$. Lemma~\ref{stabilizer} proves that this map
is injective \emph{on vertices incident with this face}: if two
representatives $h,h'\in H_i$ have the same signed vector, then
$h'^{-1}h$ belongs to
\[
         H_i\cap\operatorname{Stab}(\{v_j,-v_j\})=\{1,b\}
         \subseteq H_j.
\]
We do not identify $H_j$ with an entire torus normalizer.

Representing the vertex also by an element of $H_k$, and using
the signed stabilization of $v_i,v_k$, gives exactly the equations
of Lemma~\ref{three}. There are therefore at most three vertices
incident with both $H_i,H_k$. Three are always present:
\[
                         H_j,\quad aH_j,\quad baH_j.       \tag{2}
\]
For example $ac\in aH_j\cap H_k$ and
$bac\in baH_j\cap H_k$. They are distinct. The first is different
from the other two by $H_i\cap H_j=\gen b$. Equality of the last
two would imply $aba\in H_i\cap H_j=\gen b$, and hence $aba=b$;
this would make $a,b$ commute, contrary to hypothesis.

The only edges incident with $H_i,H_k$ are $V,bV$. To see
completeness, choose representatives $h\in H_i$, $l\in H_k$ for
any such edge. Their quotient lies in
$V=\gen a\gen{ac}$. Absorbing its two factors into $h,l$ gives
a representative in $H_i\cap H_k=\{1,b\}$.
The vertices of $V$ are $H_j,aH_j$; those of $bV$ are
$H_j,baH_j$. Thus the exact residue is the asserted path, and
every face--vertex--Petrie flag extends.

A flag containing an edge extends because its three nonedge types
form the local complete three-partite structure with two choices
per type. Any incident pair avoiding edges extends by choosing a
common group element. The only additional flags were the triples
just handled. Hence every flag extends. The eight local chambers
at an edge split into two free $V$-orbits, yielding two free global
$G$-orbits. At an endpoint of the displayed path the corresponding
three-type flag has just one completing edge, so the geometry is
not firm.
\end{proof}

\section{The required infinite Class III family}
For every odd prime power $q$, the separately reviewed cubic
generation theorem supplies a marking of $G=\PSL_2(q^3)$ of
Wilson class III. Its triality cycles $a,c,ac$ while fixing $b$,
so the three periods are equal. They exceed two: period two would
make all three generators commute, contradicting this group, and
period one is incompatible with the triality. Thus $b$ commutes
with none of the three endpoint involutions, as required above.
Theorem~\ref{completion} now gives a genuine incidence geometry.
The separately proved faithful-map symmetry theorem applies: the
edge type has two neighbors in each other type, whereas nonedge
objects have $m>2$ edges, so every correlation fixes the edge type.
Restriction to the faithful original map excludes type transpositions.
The exact type-permutation quotient is $C_3$.

\begin{corollary}
These are infinitely many pairwise nonisomorphic rank-four incidence
geometries with trialities and no dualities. Their full automorphism
group is $G$, their correlation group is $G\rtimes C_3$, and their
full automorphism group has exactly two chamber orbits.
\end{corollary}
\begin{proof}
Every type-preserving correlation restricts to an automorphism of
the original map, hence to an element of $G$. Remove that element;
the residual correlation fixes all vertices, edges and faces and
can move only Petrie objects. Conjugate it by the known triality
to obtain one which moves only vertices and fixes all edges and
faces. Its restriction to the original map is an element of $G$
fixing every edge. That edge action is faithful: its kernel is
normal and contained in $V$, whereas the standard simple group
$\PSL_2(q^3)$ has no such nontrivial normal subgroup.
The residual correlation is therefore the identity. This proves
the full automorphism group assertion. The actual order-three
triality supplies the complement to $G$ in the correlation group.

The number of edges is $|\PSL_2(q^3)|/4$, which is unbounded
and strictly increasing with $q$. Even an isomorphism allowed to
permute types must preserve the intrinsically distinguished edge
type. Hence distinct odd $q$ give nonisomorphic geometries.
The two-orbit assertion is Theorem~\ref{completion} with the now
identified full automorphism group.
\end{proof}

\paragraph{Dependencies and exact source scope.}
The entire four-page cubic-generation proof, its stated classification
and descent inputs, and the five-page faithful-map symmetry proof are
paid dependencies. Simplicity of $\PSL_2(Q)$ for $Q\geq27$ is
retained for the full-group corollary. Completion is proved by the
written algebra, with no finite premise. This recovers the odd
infinite family, not the false unrestricted Lemma 7. The final broad
question about infinite \emph{flag-transitive} families is answered by
the separately accepted later Leemans--Stokes--Tranchida work (L015),
not by this completion paper. That result is not a premise here.
No new question, Lean or numerical credit is claimed.
\end{document}
